\SourcePage{567}{570}
\section{Current density in a magnetic field}

We shall derive the quantum-mechanical expression for the current density of a
charged particle moving in a magnetic field. We start from the formula\footnote{In
this section $\mathbf j$ denotes the electric-current density: the particle-flux
density multiplied by the charge $e$.}
\begin{equation}
 \delta H=-\frac1c\int\mathbf j\,\delta\mathbf A\,dV,
 \tag{115.1}
\end{equation}
which determines the change in the Hamiltonian function of charges distributed
in space when the vector potential is varied.\footnote{The Lagrangian function
of a charge in a magnetic field contains the term
$(e/c)\mathbf v\mathbf A$, or, if the charge is represented as distributed in
space, $(1/c)\int\mathbf j\mathbf A\,dV$. Consequently, the change in the
Lagrangian function under a variation of $\mathbf A$ is
\[
 \delta L=\frac1c\int\mathbf j\,\delta\mathbf A\,dV.
\]
An infinitesimal change in the Hamiltonian function is the negative of the
change in the Lagrangian function (see I, §~40).} In quantum mechanics this
formula must be applied to the mean value of the Hamiltonian of the charged
particle:
\begin{equation}
 \overline H=\int\Psi^*\left[
 \frac1{2m}\left(\hat{\mathbf p}-\frac ec\mathbf A\right)^2
 -\frac\mu s\mathbf H\hat{\mathbf s}\right]\Psi\,dV.
 \tag{115.2}
\end{equation}
Performing the variation and recalling that
$\delta\mathbf H=\Curl{\delta\mathbf A}$, we find
\begin{align}
 \delta\overline H={}&\int\Psi^*\left[-\frac e{2mc}
 (\hat{\mathbf p}\,\delta\mathbf A+\delta\mathbf A\hat{\mathbf p})
 +\frac{e^2}{mc^2}\mathbf A\,\delta\mathbf A\right]\Psi\,dV\notag\\
 &-\frac\mu s\int\Curl{\delta\mathbf A}\cdot
 \Psi^*\hat{\mathbf s}\Psi\,dV.
 \tag{115.3}
\end{align}
We transform the term containing $\hat{\mathbf p}\,\delta\mathbf A$ by
integration by parts:
\[
 \int\Psi^*\hat{\mathbf p}\,\delta\mathbf A\Psi\,dV
 =-i\hbar\int\Psi^*\nabla(\delta\mathbf A\Psi)\,dV
 =i\hbar\int\delta\mathbf A\Psi\nabla\Psi^*\,dV
\]
(the integral over the infinitely distant surface vanishes, as usual). We also
integrate by parts in the last term of (115.3), using the familiar formula of
vector analysis
\[
 \mathbf a\cdot\Curl{\mathbf b}
 =-\Div{\mathbf a\times\mathbf b}
 +\mathbf b\cdot\Curl{\mathbf a}.
\]

\SourcePage{568}{571}
The integral of the term containing $\Div{}$ vanishes, leaving
\[
 \int\Psi^*\hat{\mathbf s}\Psi\cdot\Curl{\delta\mathbf A}\,dV
 =\int\delta\mathbf A\cdot\Curl{(\Psi^*\hat{\mathbf s}\Psi)}\,dV.
\]
The final result is therefore
\begin{align*}
 \delta\overline H={}&-\frac{ie\hbar}{2mc}\int\delta\mathbf A
 (\Psi\nabla\Psi^*-\Psi^*\nabla\Psi)\,dV\\
 &+\frac{e^2}{mc^2}\int\mathbf A\,\delta\mathbf A\Psi\Psi^*\,dV
 -\frac\mu s\int\delta\mathbf A\cdot\Curl{(\Psi^*\hat{\mathbf s}\Psi)}\,dV.
\end{align*}
Comparing this with (115.1), we obtain the following expression for the current
density:
\begin{equation}
 \mathbf j=\frac{ie\hbar}{2m}
 [ (\nabla\Psi^*)\Psi-\Psi^*\nabla\Psi]
 -\frac{e^2}{mc}\mathbf A\Psi^*\Psi
 +\frac\mu s c\Curl{(\Psi^*\hat{\mathbf s}\Psi)}.
 \tag{115.4}
\end{equation}
We emphasize that although the vector potential occurs explicitly in this
expression, the expression is nevertheless, as it must be, entirely
unambiguous. This is readily verified by direct calculation: along with the
gauge transformation of the vector potential, one must, by (111.8), also
transform the wave function according to (111.9).

It is also readily checked that the current (115.4), together with the charge
density $\rho=e|\Psi|^2$, obeys the continuity equation, as required:
\[
 \frac{\partial\rho}{\partial t}+\Div{\mathbf j}=0.
\]
The last term in (115.4) contributes to the current density through the
particle's magnetic moment. It has the form $c\Curl{\mathbf m}$, where
\begin{equation}
 \mathbf m=\frac\mu s\Psi^*\hat{\mathbf s}\Psi
 =\Psi^*\hat{\boldsymbol\mu}\Psi
 \tag{115.5}
\end{equation}
is the spatial density of magnetic moment.

Expression (115.4) is the mean value of the current density. It may be regarded
as a diagonal matrix element of an operator---the current-density operator
$\hat{\mathbf j}$. This operator is written most simply in the
second-quantized representation: one replaces $\Psi$ and $\Psi^*$ by the
operators $\hat\Psi$ and $\hat\Psi^+$ (and, by the general rule,
$\hat\Psi^+$ must stand to the left of $\hat\Psi$ in every term). The
off-diagonal matrix elements of this operator may also be defined:
\begin{equation}
 \mathbf j_{nm}=\frac{ie\hbar}{2m}
 [ (\nabla\Psi_n^*)\Psi_m-\Psi_n^*\nabla\Psi_m]
 -\frac{e^2}{mc}\mathbf A\Psi_n^*\Psi_m
 +\frac\mu s c\Curl{(\Psi_n^*\hat{\mathbf s}\Psi_m)}.
 \tag{115.6}
\end{equation}
